{{دیوان
| عنوان = ہسپانیہ
| ماخذ = ویکی ماخذ
| ماخذ_لنک = https://ur.wikisource.org/wiki/%DB%81%D8%B3%D9%BE%D8%A7%D9%86%DB%8C%DB%81
| مقابلہ = اقبال سائبر لائبریری (اقبال اکادمی پاکستان)
| مقابلہ_لنک = https://www.iqbalcyberlibrary.net
}}
=== (ہسپانیہ کی سرزمین میں لِکھّے گئے) ===

=== (واپس آتے ہوئے) ===

<poem>
ہسپانیہ تُو خُونِ مسلماں کا امیں ہے
مانندِ حرم پاک ہے تُو میری نظر میں

پوشیدہ تری خاک میں سجدوں کے نشاں ہیں
خاموش اذانیں ہیں تری بادِ سحَر میں

روشن تھیں ستاروں کی طرح ان کی سنانیں
خیمے تھے کبھی جن کے ترے کوہ و کمر میں

پھر تیرے حسینوں کو ضرورت ہے حِنا کی؟
باقی ہے ابھی رنگ مرے خُونِ جگر میں!

کیونکر خس و خاشاک سے دب جائے مسلماں
مانا، وہ تب و تاب نہیں اس کے شرر میں

غرناطہ بھی دیکھا مری آنکھوں نے و لیکن
تسکینِ مسافر نہ سفر میں نہ حضَر میں

دیکھا بھی دِکھایا بھی، سُنایا بھی سُنا بھی
ہے دل کی تسلّی نہ نظر میں، نہ خبر میں!
</poem>
